\section{为什么要把黎曼面“代数化”？}
\label{sec:ch07-motivation}
% ============================================================================

到第\ref{ch:riemann-surfaces}章为止，我们已经把 $X_0(N)$ 构造成一条紧黎曼面；
第\ref{ch:dimension}章又告诉我们，这条曲线上的微分、除子与维数公式彼此紧密相连；
第\ref{ch:jacobians}章进一步把这些数据压缩进 Jacobian 与模 Abel 簇。
从复分析角度看，这条主线已经非常完整：模曲线是复曲面上的几何对象，
权 $2$ 模形式是它的全纯微分，Hecke 算子则是微分与 Jacobian 上的自同态。

但真正的算术问题马上逼着我们再向前走一步。
Wiles 定理里出现的是“定义在 $\Q$ 上的椭圆曲线”，
不是“某个复环面”；Galois 表示里出现的是 $\Gal(\bar\Q/\Q)$ 的作用，
不是单纯的双全纯自同构群；模性定理断言的也是
“某条 $E/\Q$ 来自某个新形式”，而不是“某条复曲线长得像某个商空间”。
换句话说，如果我们只把 $X_0(N)$ 当作复解析对象，那么许多最核心的数论句子根本没有地方可落脚。

\textbf{我们遇到了什么问题？}
问题不在于“$X_0(N)$ 不是曲线”，而在于“它现在只有复解析结构”。
确实，GAGA 原理告诉我们：紧黎曼面与复射影代数曲线本质上是同一种东西。
所以从复数域观点看，把 $X_0(N)$ 看成黎曼面还是看成复代数曲线，并不会丢掉信息。
真正缺的，是\textbf{定义域}。
同一条复曲线也许可以由许多不同系数域的方程给出，
而“是否能由 $\Q$ 上的方程给出”恰恰是算术信息。
没有这层信息，我们既谈不上 $\Q$-有理点，也谈不上 Galois 共轭，更谈不上把 Hecke 算子与 Frobenius 放到同一个框架里比较。

\textbf{核心思想是什么？}
核心思想可以一句话概括：
\emph{把 $X_0(N)$ 从“复解析商空间”升级为“定义在 $\Q$ 上的代数曲线”。}
一旦做到这一点，复点集合 $X_0(N)(\C)$ 仍然是熟悉的黎曼面，
但曲线本身却拥有额外的代数结构：
它有函数域 $\Q(X_0(N))$，有 $\bar\Q$-点，
有 $\Gal(\bar\Q/\Q)$ 的自然作用，
还有在各个素数模 $p$ 下的约化。
这些正是后面 Eichler--Shimura 关系与 Galois 表示理论真正需要的语言。

\textbf{引入之后能得到什么？}
首先，$X_0(N)$ 会成为一条真正的算术曲线，
于是“点”不再只是复数参数，而可以解释成某种带附加结构的椭圆曲线。
其次，Hecke 对应会变成定义在 $\Q$ 上的代数对应，
从而不仅能作用在复微分上，也能作用在模 $p$ 约化后的几何对象上。
最后，模解释会告诉我们：
$X_0(N)$ 的 $\Q$-模型不是事后硬凑出来的，
而是由“参数化椭圆曲线及其循环子群”这个分类问题天然决定的。
这就是本章真正要完成的转变。

\begin{figure}[ht]
\centering
\begin{tikzpicture}[>=Stealth, font=\small, node distance=2.2cm]
  \node[draw, rounded corners, fill=blue!8, minimum width=4.6cm, minimum height=1cm] (riemann)
    {紧黎曼面 $X_0(N)(\C)$};
  \node[draw, rounded corners, fill=green!8, right=3.2cm of riemann, minimum width=4.8cm, minimum height=1cm] (curve)
    {复射影代数曲线 $C_{\C}\subset \mathbb P^n$};
  \draw[<->, thick] (riemann) -- node[above] {GAGA 视角} (curve);

  \node[draw, rounded corners, fill=orange!10, below=1.8cm of riemann, minimum width=4.8cm, minimum height=1cm] (qmodel)
    {$\Q$-模型 $C/\Q$};
  \node[draw, rounded corners, fill=purple!10, right=3.2cm of qmodel, minimum width=4.8cm, minimum height=1cm] (arith)
    {$C(\Q)$、$\Gal(\bar\Q/\Q)$、Hecke};

  \draw[->, thick] (curve) -- node[right] {系数域下降} (arith);
  \draw[->, thick] (qmodel) -- node[above] {复化} (curve);
  \draw[->, thick] (qmodel) -- node[above] {有理点与 Galois 作用} (arith);
\end{tikzpicture}
\caption{本章的核心转变：$X_0(N)$ 先作为复解析对象出现，但真正的算术问题要求它来自某条定义在 $\Q$ 上的代数曲线。只有在 $\Q$-模型上，Galois 作用与有理点才有内在意义。}
\label{fig:ch07-motivation}
\end{figure}


为了把上面的直觉说得精确，我们先给出本章最重要的关键词。

\begin{definition}[$\Q$-模型]
\label{def:ch07-q-model}
设 $X$ 是一条紧黎曼面。
若存在一条定义在 $\Q$ 上的光滑射影代数曲线 $C/\Q$，
使得其复化 $C_{\C}=C\times_{\operatorname{Spec}\Q}\operatorname{Spec}\C$ 的复点集合
$C(\C)$ 与 $X$ 双全纯同构，
则称 $C$ 是 $X$ 的一个\textbf{$\Q$-模型}\index{有理模型!$\Q$-模型}。
\end{definition}

这个定义看似只是多加了一层“系数落在 $\Q$ 中”的要求，
其实它正是复几何与数论之间的分水岭。
同一条复曲线只要一旦被指定为 $\Q$-模型，
就立刻得到一套之前完全不存在的算术操作。

\begin{proposition}[从 $\Q$-模型得到 Galois 作用]
\label{prop:ch07-galois-action}
设 $C/\Q$ 是一条光滑射影曲线。
则 $\Gal(\bar\Q/\Q)$ 自然作用在 $C(\bar\Q)$ 上，
并且也自然作用在函数域 $\bar\Q(C)=\bar\Q\otimes_{\Q}\Q(C)$ 上。
\end{proposition}

\begin{proof}
先看点集。
把 $C$ 嵌入某个射影空间 $\mathbb P^n_{\Q}$，设它由系数在 $\Q$ 中的齐次方程组
\[
  F_1=\cdots=F_r=0
\]
切出。
若 $P=[x_0:\cdots:x_n]\in C(\bar\Q)$，对任意 $\sigma\in\Gal(\bar\Q/\Q)$ 定义
\[
  \sigma(P)=[\sigma(x_0):\cdots:\sigma(x_n)].
\]
因为每个 $F_i$ 的系数都在 $\Q$ 中，故
\[
  F_i\bigl(\sigma(x_0),\dots,\sigma(x_n)\bigr)
  =\sigma\bigl(F_i(x_0,\dots,x_n)\bigr)=\sigma(0)=0.
\]
所以 $\sigma(P)$ 仍在 $C(\bar\Q)$ 中。
这就给出了 $C(\bar\Q)$ 上的 Galois 作用。

再看函数域。
元素 $f\in\bar\Q(C)$ 可以写成局部坐标中的有理函数，其系数落在 $\bar\Q$ 中。
定义
\[
  (\sigma f)(P)=\sigma\bigl(f(\sigma^{-1}P)\bigr).
\]
若 $f=g/h$，其中 $g,h$ 是系数在 $\bar\Q$ 中的多项式，
那么 $\sigma f$ 正是把 $g,h$ 的系数逐个施以 $\sigma$ 后所得的有理函数。
这一定义与乘法、加法相容，因此给出 $\bar\Q(C)$ 上的域自同构。
故 $\Gal(\bar\Q/\Q)$ 也自然作用在函数域上。
\end{proof}

\textbf{注。}
上面这个命题之所以重要，恰恰在于它说明：
Galois 作用不是凭空附会到复点集合上的，
而是由“方程系数落在 $\Q$ 中”这个事实强制给出的。
这也是为什么单纯知道一条紧黎曼面的复结构还远远不够。

最容易先看清的例子是级别 $1$。
在第\ref{ch:riemann-surfaces}章里，
我们已经知道 $X(1)=\SL_2(\Z)\backslash \HH^*$ 是一条亏格为 $0$ 的紧黎曼面；
而第\ref{ch:eisenstein}章构造的 $j$-函数则给出它的全局参数。
因此 $X(1)$ 不只是“某条复球面”，它实际上就是最简单的 $\Q$-曲线 $\mathbb P^1_{\Q}$。
后面我们会看到，所有 $X_0(N)$ 的 $\Q$-模型都可以看成这件事的高层级推广。

\begin{example}[为什么有理点是额外信息]
\label{ex:ch07-rational-points-extra-data}
设 $C/\Q$ 是平面圆锥曲线
\[
  x^2+y^2=z^2\subset \mathbb P^2.
\]
它的复点集合 $C(\C)$ 与 $\mathbb P^1(\C)$ 双全纯同构，
所以从复解析观点看，$C$ 与射影直线没有差别。
但从算术观点看，$C(\Q)$ 包含无穷多个勾股三元组，
而这些有理点恰恰依赖于方程系数来自 $\Q$ 这一事实。

模曲线也是同样的道理：
$X_0(N)(\C)$ 作为复曲线可以与许多别的模型同构，
但只有当我们固定了 $\Q$-模型之后，
“尖点是否有理”“CM 点如何分裂”“Hecke 对应是否与 Galois 相容”这些算术问题才变成可问、可答的问题。
\end{example}

本章接下来的安排也围绕这个目标展开：
下一节先补足最小限度的代数曲线语言；
再说明如何用 $j$-函数和模解释给出 $X_0(N)$ 的 $\Q$-模型；
随后把 Hecke 算子改写成代数对应；
最后再把整个故事从 $\Q$ 推向整数环与模 $p$ 约化。
